\subsection{nml\_time\_ave}
\label{SEC:nml_time_ave}
\inp*[nmlTimeAve]{}[nml\_time\_ave][]

This \hyperref[SEC:User_Input]{namelist block}
controls time averaging of solution variables.

\begin{nmlblocklst}[
,label={LST:nml_time_ave}
,caption={Namelist nml\_time\_ave.}
]
&nml_time_ave
  Compute_Time_Averaging      =  F
  Num_Saved_Time_Ave_Restarts =  2
  Input_Time_Ave_File         = "out.ave"
  Time_Ave_Vel_Is_Axisymm     =  F
  Start_Averaging_Timestep    = -1
  Start_Averaging_Tstar       = -1.0
/
\end{nmlblocklst}

%==========================================================================

\begin{nmldescitem}
Compute_Time_Averaging = F
\end{nmldescitem}
\inp*[nmlTimeAve::ComputeTimeAveraging]{}[Compute\_Time\_Averaging][]

\begin{nmldescbody}
This flag enables or disables the accumulation of time-averaged
flow variables.
%
If enabled, time-averaged restart and CGNS solution files will be
written whenever the standard restart and CGNS solution files are written.


%%BegCut:
%%Describe the restart files using a single paragraph with no line breaks.
%%This was hard to read, so it was replaced with the paragraph below it.
% 
% The time-averaged restart file will be written to the
%   ``\inp[nmlOutput::OutputDir]{Output\_Dir}\slash{Restart\_files}\slash''
% directory.
% %
% In order to save disk space, the name of the time-averaged restart file
% will alternate between
%   ``out.1.ave'' and
%   ``out.2.ave''.
% %
% After successfully writing a new time-averaged restart file, a
% symbolic link will be created that points to the latest
% time-averaged restart file.
% %
% The name of this link will be:
%   ``out.current.ave''.
%%EndCut:


The time-averaged restart file will be written to the
  ``\inp[nmlOutput::OutputDir]{Output\_Dir}\slash{Restart\_files}\slash''
directory.
%
In order to save disk space, the name of the time-averaged restart file
will cycle between the following:
\newline\indent
``\inp[nmlOutput::OutputDir]{Output\_Dir}\slash{Restart\_files}\slash{out.1.ave}''
\newline\indent
``\inp[nmlOutput::OutputDir]{Output\_Dir}\slash{Restart\_files}\slash{out.2.ave}''
\newline\indent
\ldots
\newline\indent
``\inp[nmlOutput::OutputDir]{Output\_Dir}\slash{Restart\_files}\slash{out.N.ave}''
\newline
where N=\inp[nmlTimeAve::NumSavedTimeAveRestarts]{Num\_Saved\_Time\_Ave\_Restarts}.
%
After a new time-averaged restart file is successfully written,
%After successfully writing a new time-averaged restart file,
a symbolic link will be created that points to the latest time-averaged restart
file. The name of this link will be:
\newline\indent
``\inp[nmlOutput::OutputDir]{Output\_Dir}\slash{Restart\_files}\slash{out.current.ave}'' 
\end{nmldescbody}

%==========================================================================

\begin{nmldescitem}
Num_Saved_Time_Ave_Restarts = 2
\end{nmldescitem}
\inp*[nmlTimeAve::NumSavedTimeAveRestarts]{}[Num\_Saved\_Time\_Ave\_Restarts][]

\begin{nmldescbody}
This is used to change the number of time-averaged restart files
(out.$\ast$.ave) that are saved to the Restart\_files directory.
The default is 2, so output will cycle between out.1.ave and out.2.ave.

Through the setting of
\inp[nmlOutput::RestartInterval]{Restart\_Interval},
\inp[nmlOutput::RestartFrequency]{Restart\_Frequency}, and
this input,
one should ensure there is a viable fallback solution to
restart from should an NaN occur.
\end{nmldescbody}

%==========================================================================

%##%\FINALFIX{\newpage}
\begin{nmldescitem}
Input_Time_Ave_File = "out.ave"
\end{nmldescitem}
\inp*[nmlTimeAve::InputTimeAveFile]{}[Input\_Time\_Ave\_File][]

\begin{nmldescbody}
\IDEA{Change this input to \inptext{Input\_Time\_Ave\_File}.}

File path to the time-averaged solution restart file that will be
read during initialization if the following criteria are met:
%
\begin{itemize}
  \item
    $\inp[nmlInitialization::LoadRestartFile]{Load\_Restart\_File}=\mathtt{T}$
  \item One of the following is met:
    \begin{itemize}
      \item
        $\inp[nmlTimeAve::ComputeTimeAveraging]{Compute\_Time\_Averaging}=\texttt{T}$
      \item
        $\inp[nmlTimeAve::StartAveragingTstar]{Start\_Averaging\_Tstar}>0$ and
        $\inptext{Start\_Averaging\_Tstar} \leq \tstar$,
        where \key[key:tstar]{\tstar}[t *@\tstar] is the
        non-dimensional simulation time, \eref{EQN:tstar}, of the data that was
        read in from \inp[nmlInitialization::RestartFile]{Restart\_File}.
      \item
        $\inp[nmlTimeAve::StartAveragingTimestep]{Start\_Averaging\_Timestep}>0$
        and $\inptext{Start\_Averaging\_Timestep} \leq iteration$,
        where $iteration$ is the iteration number of the data that was
        read in from \inptext{Restart\_File}.
    \end{itemize}
\end{itemize}
%
The code will read the time-averaged flow variables of the
previous simulation from which the new simulation was restarted,
and then add to those statistics during the current run job.

\NOTE
If the above criteria are met, but the file specified by
\inptext{Input\_Time\_Ave\_File} does not exist, a completely new time averaging will
be started.
%
This is the case even if you were intending to restart a previous time
averaging, and the reason the file could not be found was due to a small typo in
the file name.

\NOTE
This file is only for \emph{reading} during initialization and
is \emph{not an output file}.
\end{nmldescbody}

%==========================================================================

\begin{nmldescitem}
Time_Ave_Vel_Is_Axisymm = F
\end{nmldescitem}
\inp*[nmlTimeAve::TimeAveVelIsAxisymm]{}[Time\_Ave\_Vel\_Is\_Axisymm][]

\begin{nmldescbody}
% Specifies whether axisymmetric velocity components should be included
% in the time-averaged flow variables.
% 
% If false, the time-averaged Cartesian velocity components
% ($v_x$, $v_y$, $v_z$) will be saved
% along with the Cartesian stress and heat flux data.
% %
% If true, the grid is assumed to be axisymmetric about the $x$-axis and
% the time-averaged
%   streamwise      ($v_x$),
%   radial          ($v_r$), and
%   circumferential ($v_\theta$)
% velocity components will added.
% %
% This option is useful for recording statistics for axisymmetric flows.

This option is useful for recording statistics for axisymmetric flows.
%
It specifies whether axisymmetric velocity components should be included
in the time-averaged flow variables.
%
Note, regardless of the value for this input variable, the Cartesian momentum
and velocity components ($[\rho v_x,\rho v_y,\rho v_z]$ and $[v_x,v_y,v_z]$)
will always be averaged and saved to the time-averaged restart file.
\\[0.25em]
\begin{tabular}{ @{\qquad} r @{ \quad } p{0.85\linewidth} @{} }
  T & The geometry is assumed to be axisymmetric about the $x$-axis.
      %
      The axisymmetric momentum and velocity components in the
      radial ($\rho v_r, v_r$) and
      circumferential ($\rho v_\theta, v_\theta$)
      directions will also be averaged and saved to the
      time-averaged restart file.
      %
      The stress and heat flux data will be for the axisymmetric terms,
      \emph{not} the Cartesian terms
      (\thatis the axisymmetric velocity components are used in the
      time-averaged terms).
      %
      For example, the averaged term \ReyAve{v_y T} will be replaced with
      \ReyAve{v_r T} and \ReyAve{\rho v_x v_z} will be replaced with
      \ReyAve{\rho v_x v_\theta}.
  \\
  F & % The time-averaged Cartesian velocity components
      % ($v_x$, $v_y$, $v_z$) will be saved.
      %
      The stress and heat flux data will be for the Cartesian terms
      (\thatis the Cartesian velocity components are used in the
      time-averaged terms).
\end{tabular}
\\[0.25em]
\NOTE
See \sref{SEC:Files:Time_Averaging_Output} for additional details.

\end{nmldescbody}

%==========================================================================

\begin{nmldescitem}
Start_Averaging_Timestep = -1
\end{nmldescitem}
\inp*[nmlTimeAve::StartAveragingTimestep]{}[Start\_Averaging\_Timestep][]

\begin{nmldescbody}
Iteration number to begin time averaging.
\end{nmldescbody}

%==========================================================================

\begin{nmldescitem}
Start_Averaging_Tstar = -1.0
\end{nmldescitem}
\inp*[nmlTimeAve::StartAveragingTstar]{}[Start\_Averaging\_Tstar][]

\begin{nmldescbody}
Non-dimensional simulation time (\key[key:tstar]{\tstar}[t *@\tstar]),
\eref{EQN:tstar},
to begin time averaging.
\end{nmldescbody}

%==========================================================================

